% Generated draft: requires mathematical and source review.
\documentclass[11pt]{article}
\usepackage[margin=25mm]{geometry}
\usepackage{fontspec}
\setmainfont{Latin Modern Roman}
\usepackage{amsmath,amssymb,mathtools,unicode-math}
\setmathfont{Latin Modern Math}
\usepackage{enumitem,needspace,xurl,hyperref}
\setlength{\parindent}{0pt}
\setlength{\parskip}{5pt}
\setlength{\emergencystretch}{3em}
\newcommand{\N}{\mathbb{N}}
\newcommand{\Z}{\mathbb{Z}}
\newcommand{\Q}{\mathbb{Q}}
\newcommand{\R}{\mathbb{R}}
\newcommand{\C}{\mathbb{C}}
\newcommand{\F}{\mathbb{F}}
\newcommand{\A}{\mathbb{A}}
\newcommand{\PP}{\mathbb{P}}
\newcommand{\E}{\mathbb{E}}
\newcommand{\eps}{\varepsilon}
\newcommand{\dd}{\,\mathrm{d}}
\begin{document}
\section*{UnsolvedMath ID: 20001886}
\textbf{Source problem code:} catalogue entry
\subsection{Definitions and mathematical statement}
Schmidt's source asks whether the marked length spectrum determines a locally \(\operatorname{CAT}(-1)\) metric on a two-dimensional simplicial complex, either on the fixed space or among homeomorphic complexes with the preceding properties. Those preceding properties are absent from this record. With a fixed marking of fundamental groups, the marked length function assigns to a conjugacy class the infimum of lengths of loops in that class. Equality means equality on every marked class, not merely equality of the unmarked set of lengths. The requested equivalence is a marking-preserving isometry.

The source attributes a positive topological-manifold case to [Her-Pau1997], asks also about higher dimensions, and singles out the thick case where each edge belongs to at least three two-simplices. The aid reports later rigidity for compact Fuchsian-building quotients and simple thick surface amalgams while the arbitrary locally \(\operatorname{CAT}(-1)\) two-complex question remains open in its checked literature. Exact hypotheses and references for these subclasses are not supplied; thickness alone is not asserted to settle the whole question.

The candidate finite criterion fixes one finite connected triangulation \(K\), with labelled edges \(e_1,\ldots,e_m\), and two locally \(\operatorname{CAT}(-1)\) piecewise-hyperbolic metrics \(g,g'\) whose nondegenerate faces are hyperbolic triangles of curvature \(-1\). Suppose there are closed edge circuits \(\gamma_1,\ldots,\gamma_s\) that are locally geodesic for both metrics. Local geodesicity at vertices is a metric link-angle condition, not guaranteed by using the same combinatorial word. Define the unoriented traversal matrix
\[
 A_{ij}=\#\{\text{traversals of edge }e_j\text{ by }\gamma_i\},
 \qquad A\in M_{s\times m}(\Z_{\ge0}).
\]
Assume \(\operatorname{rank}A=m\). The aid reports that their marked lengths then give
\[
 (\ell_g(\gamma_i))_{i=1}^s=A(\ell_g(e_j))_{j=1}^m,
\]
and the identical relation for \(g'\). Local geodesicity ensures these circuit lengths are the class's marked geodesic lengths in the compact locally negative-curvature setting. Equal values for these finitely many classes and full column rank recover every labelled edge length equally. Hyperbolic side-side-side congruence then gives compatible face isometries and a marking-preserving isometry of the complete labelled simplicial metrics.

This is a sufficient calibration condition, not a proof that every complex has such circuits, not a necessary condition for all marked-length rigidity, and not a theorem for variable-curvature face metrics or different unaligned triangulations. Signed homology incidence cannot replace the stated unoriented traversal counts.

The adjacent boundary-rigidity discussion considers a compact manifold \(M\) with boundary \(N\) and its boundary distance function \(d_g(p,q)\), \(p,q\in N\). A meaningful usual question is whether equality with \(d_{g'}\) forces an isometry between \((M,g)\) and \((M,g')\), fixing the boundary marking. The extraction instead ends by saying the new metric must be isometric to itself; that tautological wording is explicitly not taken as a theorem. Boundary data and marked closed-orbit data are related motivations, not identified by an unprovided argument here.
\subsection{Short English statement}
The aid gives a finite sufficient test on one fixed hyperbolic triangulation: common locally geodesic edge circuits with full-rank unoriented traversal matrix recover all edge lengths, then all triangle metrics. It does not establish such a calibration for every locally CAT(-1) complex or settle the broader rigidity question.
\subsection{Sources}
\begin{itemize}
\item Local UnsolvedMath registry, permanent ID \(20001886\), source code \texttt{catalogue entry}; \texttt{aim-geometry-notes.json}, \texttt{notes[223]}, tag \texttt{question}.
\item AIM workshop ``Geodesics,'' recorded item 5.1, extracted Question 5.1.2 (Schmidt): \url{https://aimath.org/pastworkshops/geodesicsproblems.pdf}. Source type PDF; statement status \texttt{exact}. Supplied verification identifies the record. Only the supplied text was used; no PDF was accessed or compiled.
\item Source reference label [Her-Pau1997] for the manifold case. The aid names compact Fuchsian-building quotients and simple thick surface amalgams without authors, titles, dates, identifiers, URLs, or authoritative full citations. The missing preceding properties and finite-criterion proof artifacts are not supplied; original remarks and literature fields are empty.
\item Research attempt 1: \texttt{partial\_result}, status at run \texttt{partially\_solved}, rigor label \texttt{proved}, novelty confidence low; registry status \texttt{partially\_solved}, classification \texttt{PARTIAL-PROGRESS}. Created and updated 14 August 2026. Machine-generated aid, not peer reviewed; independently verify before citation.
\item Ulam AI / UnsolvedMath Contributors; CC BY 4.0.
\end{itemize}
\end{document}
